% 2026-11-23 站主問：「兩人彼此有什麼的誤解」「嚴重的認知落差」。
% 兩篇文字層主對話讀過；Post 註 33、註 45 註號依 JSTOR 本文字層（註號未拆散）。誰誤讀誰為整理者依原文所作之判斷，正文逐處寫出所據頁碼。
\chapter{兩文之互相誤讀}
\label{ch:70_misreadings}

\begin{qaquestion}
Michelman\ix{persons}{michelman@Michelman, Frank I.} 與 Post\ix{persons}{post@Post, Robert} 兩文互為回應，彼此有哪些誤讀？
\begin{enumerate}
\tightlist
\item Michelman 如何描述 Post 的民主觀，Post 如何回應？
\item Post 如何重述 Michelman 的模型，與 Michelman 原文有何出入？
\item 兩人實際一致之處何在？
\end{enumerate}
\end{qaquestion}

Post 的評論依次回應 Michelman 第 III、V、VII 部。兩文對三件事的說法互有出入：民主屬程序概念或實質概念；Michelman 第 VII 部的模型要求公民同意判決結果，或要求公民尊重作成判決的程序；Post 所駁斥的立場，Michelman 本人是否接受。

\section{Michelman 將 Post 歸為純粹程序論}
\label{sec:70_misreadings-1}

Michelman 405 頁稱 Post 的回應式民主\ix{terms}{01@\zhlidxcat{民主理論}!回應式民主}「bottomlessly procedural, resting on no substantive foundation and implying no substantive presupposition」〔程序到底，不以任何實質為基礎，也不蘊含任何實質預設〕，引的是 Post 書中轉引 Lefort\ix{persons}{lefort@Lefort, Claude} 的一句：民主政體建立在一場正當與不正當之辯論的正當性上，這場辯論「necessarily without any guarantor and without any end」〔必然沒有擔保者，也沒有終點〕。\footcite[405 \& n.18]{michelman1998}407 頁 Michelman 接著指出，Post 的理論以「absolutely unrestricted access to public discourse for every person and every view」〔每個人、每種觀點均得不受限制地進入公共論述〕為前提，這個原則雖以程序事項為對象，本身卻是實質，持此說者須將它置於一切辯論之外。\footcite[407]{michelman1998}第 2 部整部以此為據，指 Post 的程序論在根基處含有實質。

Post 438 頁註 33 回應：「The procedural account of democracy that Michelman attributes to me seems roughly to express merely one implication of the substantive idea of democracy for the regulation of public discourse」〔Michelman 歸於我的程序民主觀，大致只表達了實質民主理念對公共論述之規制所生的一項推論〕。\footcite[438 n.33]{post1998}同頁正文說，人民主權\ix{terms}{01@\zhlidxcat{民主理論}!人民主權}依程序標準界定，民主依實質價值界定，該價值為集體自治。\footcite[438]{post1998}依 Post 自述，Michelman 在第 2 部揭出的實質前提，Post 原本即予主張；Michelman 該部的批評，以 Post 未採的立場為對象。Michelman 405 頁所據，僅為 Post 書中轉引 Lefort 的一段，這段話也可讀為對公共論述之開放性的描述（第 \ref{ch:20_responsive} 部）。

\section{Post 將 Michelman 的模型重述為以同意判決結果為前提}
\label{sec:70_misreadings-2}

Post 把 Michelman 425--426 頁的兩項條件命名為「Michelman's model」，並稱該模型的一項前提是「a Justice's judgments are generally accepted as `just, or fair,' or `right'」〔大法官的判斷普遍被接受為正義、公平或正確〕。\footcite[432]{post1998}433 頁他將此比於 Rousseau\ix{persons}{rousseau@Rousseau, Jean-Jacques} 使個別意志與普遍意志合一，436 頁推出：在這個模型裡，公民只在同意法院憲法判斷的範圍內才算自治，「Consensus defines and exhausts the scope of democratic self-governance」〔共識界定並窮盡民主自治的範圍〕，分歧因此在一切情形均否定民主正當性。\footcite[433, 436]{post1998}

Michelman 424--426 頁的推論分兩步。第一步，相信根本法有程序獨立之正確標準的人，所能尊重的對象為法律之正確，與他對法律內容保有的些微影響力無關。\footcite[425]{michelman1998}第二步以「Now complicate the situation」起頭：此人不認為任何單獨判斷的人可靠，自己與 Brennan\ix{persons}{brennan@Brennan, William J., Jr.} 皆然；他的尊重因此改繫於兩項條件，一為持續盡最大可行努力使根本法及其主要解釋正確，二為使有權決定者暴露於社會上各種意見與利益陳述的全面衝擊之下。\footcite[425]{michelman1998}原文第 VI 部的論證，用意相同：合理的詮釋多元使任何人無從確知何者正確，公民支持政權之正當性，只能改求於決定機關對公眾意見之開放（第 \ref{ch:40_paradox} 部）。\footcite[422--423]{michelman1998}第二步所處理的，正是不同意結果的人。Post 432 頁的重述停在第一步。

Post 436 頁另設一例：九位 McReynolds\ix{persons}{mcreynolds@McReynolds, James C.} 組成的法院，勤勉、稱職，也暴露於各種意見之下，判決卻與你的法學見解深相牴觸，Post 懷疑讀者會把這樣的法院當作自治的工具。\footcite[435--436]{post1998}此例顯示兩項條件不足以使人自治，Post 由此推出「At most, these conditions can establish only political legitimacy」〔至多，這些條件只能確立政治正當性〕。\footcite[436]{post1998}結論成立；推論所據的前提，則與 Michelman 原文第二步不合。

\section{Post 稱 Michelman 的模型排除參與}
\label{sec:70_misreadings-3}

Post 434 頁說，主張以公眾意見之形成為民主自治所在的論者，強調公民參與其所受治之法律的形成，Michelman 的模型則「explicitly excludes this notion of participation」〔明白排除這種參與概念〕。\footcite[434]{post1998}434--435 頁再以 Fuller\ix{persons}{fuller@Fuller, Lon L.} 所釋的對抗制相比：訴訟當事人受法院管轄，卻不能說他們透過法院之管轄治理自己。\footcite[434--435]{post1998}

Michelman 425 頁排除的對象，為參與作為控制力：意見與利益陳述得以登錄，要緊之處在其「presumptive epistemic value」〔推定的認識價值〕，而不在「some bit of leverage you exercise over what gets decided」〔對所作決定施加的一點槓桿力〕。\footcite[425]{michelman1998}參與本身仍在模型之內。426 頁他說 Brennan 只在法院行使權力以保護並擴張每個人向法院及其他社會權威陳述己見之權利與機會的範圍內，才承認法院有此權力；424 頁要求決定機關始終受「a public-opinion-in-formation that is bent on democratizing itself and the legal and social conditions of its production」〔致力於使自身及其生產之法律與社會條件民主化的形成中公眾意見〕之壓力。\footcite[424, 426]{michelman1998}兩人所稱的參與，Post 指參與法律之形成，Michelman 指對公共討論的認識上貢獻。

\section{Post 所駁斥的立場與 Michelman 本人的結論}
\label{sec:70_misreadings-4}

Post 432 頁承認，看不出 Michelman 本人是否接受 425--426 頁的推測，並說「there is some textual evidence that Michelman is not ultimately satisfied with it」〔有若干文本跡象顯示 Michelman 終究不滿意於此〕，仍主張該立場「ought to be definitively repudiated」〔應予明確駁斥〕。\footcite[432]{post1998}

Michelman 427 頁已寫道：「To press your views upon ruling authorities is not yet to rule. To find the laws deserving of your respect is not yet to decide the laws」〔向統治機關施加你的意見，尚不是統治。認為法律值得你尊重，尚不是決定法律。〕，並借 Churchill 之語稱 Brennan 式民主或許是除其他一切形式以外最壞的民主。\footcite[427]{michelman1998}Post 439 頁引這一句作為自己原文第 V 部的結論，442 頁對 Michelman 426 頁評 Brennan 的一段說「I agree entirely」，並稱 Brennan 的成就不能稱為「self-government through democracy」。\footcite[439, 442]{post1998}兩人對 Brennan 的判斷一致：Brennan 使法院成為值得尊重的權威，但他的判決未使不同意者因此自治。Post 第 II 至 V 部大半篇幅所駁斥的，是兩人均不接受的立場。

\section{民主是否被排除於根本法之決定以外}
\label{sec:70_misreadings-5}

Michelman 420 頁說，有一項邏輯使民主「entirely bars」〔完全被排除〕於決定一國根本法內容的空間以外。\footcite[420]{michelman1998}Post 430--431 頁以依美國聯邦憲法第 5 條\ix{statutes}{02@美國!01@聯邦憲法!005@第 5 條}修憲廢除增修條文第 26 條\ix{statutes}{02@美國!01@聯邦憲法!126@增修條文第 26 條}、規定未滿 21 歲者不得投票為例：法官即使認為修正案反民主，仍會依修正案判決，此時「popular sovereignty manifestly trumps constitutionalism」〔人民主權明顯勝過憲政主義〕；他以此指 Michelman 的結論未能說明實務上的典型案例。\footcite[430--431]{post1998}

Michelman 419 頁的論點以證立為對象：任何人都可主張現行憲法經由民主的歷史產生，但不能以這段歷史作為要求尊重現行憲法的依據，除非先預設凡經此種程序產生的憲法均應施行，而這項預設本身若再以民主歷史為據，即陷入惡性的無窮後退。\footcite[419]{michelman1998}修正案在實證法上的拘束力，Michelman 並未否認。依 Post 自己 441 頁註 45 的表，第 5 條修憲程序與民主理念之關係屬第 2 格（第 \ref{ch:60_post} 部表 6.1），Michelman 所論為規範上的證立。Post 431 頁註 10 也承認此例「loaded the dice」〔在骰子上動手腳〕：換成授予總統終身任期的修正案，民主與人民主權之衝突即告顯露。\footcite[431 n.10]{post1998}Post 437--438 頁隨後把憲政主義與民主之衝突改寫為人民主權與民主理念之衝突，並引 Michelman 419 頁「Democracy is a demanding normative idea, an idea with content」〔民主是一個要求甚高的規範理念，一個有內容的理念〕為據。\footcite[438]{post1998}這項改寫與 Michelman 419 頁的論點一致，兩人在此分歧的是用語：Michelman 以民主稱民主程序，Post 將民主程序稱為人民主權，民主一詞保留給實質價值。
